\begin{table}[htbp]
\centering
\footnotesize
\caption{Decomposition of the intercept and the plateau versus power-law comparison.}
\label{tab:s2}
\begin{tabular}{lllrrrrrr}
\toprule
Specification & $a$ total & $a$ | week FE & $a$ | week + state-season & Week share & ICC & Power exp. & $p$ & CV folds \\
\midrule
$R_A$, rule windows ($S=4$) & 0.0150 [0.0100, 0.0208] & 0.0088 [0.0064, 0.0124] & 0.0083 [0.0064, 0.0141] & 0.41 & 0.25 & 0.38 & 0.11 & 26/51 \\
$G$, rule windows ($S=4$) & 0.0840 [0.0508, 0.1226] & 0.0415 [0.0306, 0.0586] & 0.0394 [0.0293, 0.0614] & 0.51 & 0.31 & 0.35 & 0.44 & 19/51 \\
$G$, calendar windows ($S=3$) & 0.0991 [0.0583, 0.1399] & 0.0473 [0.0337, 0.0676] & 0.0450 [0.0335, 0.0932] & 0.52 & 0.29 & 0.27 & 0.42 & 28/51 \\
\bottomrule
\end{tabular}
\par\smallskip
\begin{minipage}{0.95\linewidth}\scriptsize Intervals: two-way bootstrap with $B=1000$ (separate draws from Table S1, hence small differences). Week FE: national week fixed effects removed from the statistic before the fit; state-season effects removed in addition in the next column. Week share = 1 - $a$|week FE / $a$ total. ICC: share of the statistic's variance between weeks. Power exp.: exponent $c$ of $L = \alpha D^{-c}$. $p$: parametric-bootstrap test of the affine plateau against the power law (200 null draws). CV folds: leave-one-state-out folds in which the affine model predicts better. A power analysis found no fraction of the estimated $a$ ($\le 1$) detectable with 80\% power.\end{minipage}
\end{table}
